State Ito's Formula and some general
facts about stochastic integrals


Let \( f(t, x) \) be a cont. function with cont. partial derivatives
\( \frac{df}{dt}, \frac{df}{dx}, \frac{d^{2}f}{dx^2} \), then
\( f(t, W_t) = f(a, W(a)) + \int_{a}^{t} \frac{df}{dx}(s, W(s) dW_s + \int_{a}^{t} \frac{df}{dt}(s, W_s) + \frac{1}{2} \frac{d^{2}f}{dx^2}(s, W_s) ds\)

Additionally we have Ito's product formula (requires \( X_t, Y_t \) to be diffusions/Ito processes)

\( X_tY_t = X_aY_a + \int_{a}^{t} Y_s dX_s + \int_{a}^{t} X_s dY_s + \int_{a}^{t} dX_s dY_s \)

(Not important in the lecture but) more generally for cadlag, square-integrable martingales \( M_t \) 
e.g. compensated Poisson processes, we have for 

\((f(M_t) = f(M_a) + \int_{a}^{t} F'(M_{s-}) dM_s + \frac{1}{2}\int_{a}^{t}F''(M_s) d[M]^{c}_s + \sum_{a < s \leq t} (F(M_s) - F(M_{s-}) - F'(M_s-)\Delta M_s\)
where \( [M]^c_s \) is the cont. part of \( [M]_s \) and \( \Delta M_s = M_s - M_{s-} \) and \( [M_t] = \lim_{\|\delta_n\| \to 0} \sum_{i}^{n} (M_{t_i} - M_{t_{i-1}})^2 \)
is the quadratic variation process (for cont. Martingales, this is the variance).

Let \( f \in L^2[a, b] \). The wiener integral of \( f \), \( \int{a}^{b} f(t) d_Wt \) is a Gaussian random variable
with mean \( 0 \) and variance \( \|f\|^2 = \int_{a}^{b} f^2 dt \).
The last part is a special case of Ito's isometry, which states
\( \mathbb{E}[(\int_{0}^{T} X_t dW_t)^2] = \mathbb{E}[\int_{0}^{T} X_t^2 dt] \)
for stochastic processes on the same space \( \Omega \) defined as \( W_t \)
and the probability measure for the expectation is the Wiener measure.


Elaborate on 

How can a noncentral \( \chi^2 \)-distribution be obtained from a Poisson-distribution and
\( \chi^2 \)-distributions? How can this be used for an exact simulation of a square-root
diffusion?

This also answers
More generally, how can a noncentral \( \chi^2 \) -distribution be constructed from a \( \chi^2 \)-
distribution and a normal? How can this be used for an exact simulation of a
square-root diffusion?

How is a \( \chi^2 \)-distribution related to normals if its degree of freedom is a natural
number? How can a noncentral \( \chi^2 \)-distribution be constructed similarly?

Feller \& Cox determined that the conditional distribution of square-root diffusion \( r(t) \)
given \( r(u) \) is a non-central chi-square distribution with a certain set of parameters (which are quite complex and not requested).

A chi-square distr. is typically defined as the distribution
of the following r.v.
\( Q = \sum_{i=1}^{k} Z_i^2 \), where \( Z_i \sim \mathcal{N}(0, 1) \) are i.i.d.
and \( k \) are the referred to as the degree of freedoms.

In turn a non-central \( \chi^2 \)-distr. with parameter \( \lambda \) is defined similarly but this time \( Z_i \sim \mathcal{N}(\mu_i, 0) \)
and \( \lambda = \sum_{i=1}^k \mu_i^2 \).

Its distribution is given by the rather complex term

\( \mathbb{P}[\chi^2_{k}(\lambda) \leq y = e^{\frac{\lambda}{2}}b\sum_{j=0}^{\infty} \frac{\frac{1}{j!}(\frac{1}{2}\lambda)^j}{2^{\frac{k}{2}+j}\Gamma(\frac{k}{2}+j)} \int_{0}^{y} z^{\frac{k}{2} + j -1} e^{-\frac{z}{2}}dz \)

Its simulation however is not as bad as it looks:
For example it holds for \( k > 1 \)

\( \chi^2_k(\lambda) = (Z + \sqrt{\lambda})^2 + \chi^2_{k-1} \)
so it suffices to simulate central \( \chi^2 \) and normal distr.

Another way is to use a Poisson distribution, which still works for \( k \leq 1 \)
namely to simulate first a Poisson variable \( N \) with parameter \( \frac{\lambda}{2} \)
and then a central \( \chi^2 \)-square with \( k + 2N \) degrees of freedom.

In tandem these 2 methods allow for a complete simulation of a square-root process.

Last the gamma distribution can be leveraged for the central chi-square distr., since \( X \sim \Gamma(\frac{k}{2}, 2) \iff X \sim \chi^2{k}\),
which enjoys a quite symmetric structure.
In particular the pdf of the Gamma distr. \( \Gamma(\alpha, \beta) \) is
\( \frac{1}{\Gamma(\alpha)\beta^k}x^{\alpha-1}e^{-\frac{x}{\beta}} \)
with mean \( \alpha\beta \) and variance \( \alpha\beta^2 \) and for a sum of independent \( X_i \sim \Gamma(k_i, \beta)\)
we have \( \sum_{i=1}^{N} X_i \sim \Gamma(\sum_i k_i, \beta) \) and multiplying some \( X \sim \Gamma(\alpha, \beta) \) with a constant \( c > 0 \)
gives \( cX \sim \Gamma(\alpha, c\beta) \), so it suffices to be able to simulate \( \Gamma(\alpha, 1) \) for \( \alpha \in (0, 1] \) 
(in the lecture we also consider this for \( a > 1 \), but this complicates things unnecessarily).
additionally it is enough to only simulate \( \Gamma \)


Correction to (t,s)-sequences
The set \( \{x_{kb^m}, x_{kb^m + 1}, \dots, x_{(k+1)b^{m} - 1}\}\) is a \( (t, m, s) \)-net in base \( b \)

van der Corput sequence
The \(b\)-ary representation of the positive integer \(n \ge 1\) is
\( n = \sum_{k=0}^{L-1} d_k(n) b^k = d_0(n) b^0 + \dots + d_{L-1}(n) b^{L-1} \);
where \(b\) is the base in which the number \(n\) is represented, and \(0 \le d_k(n) < b\); that is, the \(k\)-th digit in the \(b\)-ary expansion of \(n\). The \(n\)-th number in the van der Corput sequence is
\( g_b(n) = \sum_{k=0}^{L-1} d_k(n) b^{-k-1} = d_0(n) b^{-1} + \dots + d_{L-1}(n) b^{-L} \).

E.g. for \( b = 10 \) we have \( 1 = 1 \cdot 10^0,\ 2 = 2 \ \dot{c} 10^{0}, \dots, 11 = 1 \cdot 10^0 + 1 \cdot 10^1, \dots \) thus
the van der Corput sequence reads
\( \frac{1}{10}, \frac{2}{10}, \dots, \frac{1}{10} + \frac{1}{100} = \frac{11}{100}, \dots \).


Hardy-Krause variation depends only on the function which is chosen as the inte-
grand; just take it as given. State the Koksma-Hlawka bound. Why does it imply
a better rate for QMC in comparison to MC, if the star discrepancy of the sequence
of used points converges to zero at the best conjectured rate?

We won't go into detail, defining the Hardy-Krause variation \( V(f) \)
of a function \( f \), because its rather awful.
The Koksma-Hlawka bound exists if \( V(f) < \infty \)
and states
\( |\frac{1}{n}\sum_{i=1}^n f(x_i) - \int_{[0, 1]^d} f(u) du| \leq V(f)D^{\ast}(x_1, \dots, x_n) \)
while according to the central limit theorem for \( U_i \sim \text{unif}[0, 1]^{d} \) and \( \sigma^{2}_f = \text{Var}(f(U)) \)
we have
\( |\frac{1}{n}\sum_{i=1}^{n} f(U_i) - \int_{[0, 1]^d} f(u)du| \leq z_{\delta/2} \frac{\sigma_f}{\sqrt{n} } \)
with probability \( 1 - \delta \) and \( -z_{\frac{\deta}{2}} \) is the \( \frac{\delta}{2} \)-quantile of the normal distribution.

While the Koksma-Hlawka is inpractiable due to the complexity of determining \( V(f) \), which is usually straightforward for \( \sigma^{2}_f \)
It provides theoretically a better bound if \( D^{\ast}(x_1, x_2, \dots, x_n) \to 0 \) and this happens faster than \( \mathcal{O}(\frac{1}{\sqrt{n}} \)
which is at least the case for \( \mathcal{O}(\frac{\log(n)^{d}}{n}), \ d \geq 3 \) (which we determined earlier for \( D* \)).

